\exercisegroup

\begin{enumerate}
\def\labelenumi{\arabic{enumi})}
\tightlist
\item
  Let \(A\) be an infinite set.
\end{enumerate}

\begin{enumerate}
\def\labelenumi{\alph{enumi})}
\tightlist
\item
  Let E be the Banach space \(\overline{\mathcal{K}(\mathbf{A})}\) over \(\mathbf{R}\) , consisting of all families \(x = (x_{\alpha})_{\alpha \in \mathbf{A}}\) of real numbers such that \(\alpha \mapsto x_{\alpha}\) tends to 0 with respect to the filter of complements of finite subsets of A, and endowed with the norm \(\| x\| = \sup_{\alpha \in \mathbf{A}}|x_{\alpha}|\) (when \(\mathbf{A} = \mathbf{N}\) , this space is denoted by \(c_0\) or \(c_0(\mathbf{N})\) . Show that every continuous linear form on E can be written in a unique way as \(x \mapsto \sum_{\alpha \in \mathbf{A}} u_{\alpha} x_{\alpha}\) , where \((u_{\alpha})_{\alpha \in \mathbf{A}}\) is a family such that \(\sum_{\alpha \in \mathbf{A}} |u_{\alpha}| < +\infty\) ; then the dual \(\mathrm{E}'\) of E can
\end{enumerate}

be identified (as a non topological vector space) with the space \(\ell^1 (\mathbf{A})\) (I, p.~4, Example). \(b)\) Let \(\mathrm{F}\) be the Banach space \(\ell^1 (\mathbf{A})\) (I, p.~4, Example) (when \(\mathbf{A} = \mathbf{N}\) , also denoted by \(\ell^1\) ). Show that every continuous linear form on \(\mathrm{F}\) can be written in a unique way as \(x\mapsto \sum_{\alpha \in \mathbf{A}}u_{\alpha}x_{\alpha}\) , where \((u_{\alpha})_{\alpha \in \mathbf{A}}\) is a bounded family of real numbers; then the dual \(\mathrm{F}'\) of \(\mathrm{F}\) can be identified (as a non topological vector space) with the space \(\mathcal{B}(\mathbf{A}) = \ell^{\infty}(\mathbf{A})\) .

\begin{enumerate}
\def\labelenumi{\alph{enumi})}
\setcounter{enumi}{2}
\tightlist
\item
  Let B be an arbitrary set, \((c_{\alpha \beta})_{(\alpha ,\beta)\in \mathrm{A}\times \mathrm{B}}\) an arbitrary family of numbers \(\geqslant 0\) . Let G be the vector space of all families \(x = (x_{\alpha})_{\alpha \in \mathbf{A}}\) of real numbers such that, for every \(\beta \in \mathbf{B}\) , we have \(p_{\beta}(x) = \sum_{\alpha \in \mathbf{A}}c_{\alpha \beta}|x_{\alpha}| < + \infty\) ; the \(p_{\beta}\) are semi-norms on G. In order that G, with the topology defined by this family of semi-norms, be Hausdorff, it is necessary and sufficient that, for every \(\alpha \in \mathbf{A}\) , there exists at least one \(\beta \in \mathbf{B}\) such that \(c_{\alpha \beta} > 0\) . Show that then \(\mathbf{G}\) is complete, and that every continuous linear form on \(\mathbf{G}\) can be written uniquely as \(x \mapsto \sum_{\alpha \in \mathbf{A}} u_{\alpha} x_{\alpha}\) , where
\end{enumerate}

\((u_{\alpha})_{\alpha \in \mathbf{A}}\) is a family of real numbers satisfying the following condition; there exists a finite number of indices \(\beta_{i}\in \mathrm{B}(1\leqslant i\leqslant n)\) and a number \(a > 0\) such that \(|u_{\alpha}|\leqslant a.c_{\alpha \beta_i}\) for all \(\alpha \in \mathbf{A}\) and \(1\leqslant i\leqslant n\) . Prove the converse, and extend these results to the case where the field of scalars is \(\mathbf{C}\) .

\begin{enumerate}
\def\labelenumi{\arabic{enumi})}
\setcounter{enumi}{1}
\tightlist
\item
  \begin{enumerate}
  \def\labelenumii{\alph{enumii})}
  \tightlist
  \item
    Let F and G be two vector spaces in separating duality. Show that, if F is relatively bounded for \(\sigma(\mathrm{F}, \mathrm{G})\) (III, p.~43, exerc. 6), then G is relatively bounded for \(\sigma(\mathrm{G}, \mathrm{F})\) .
  \end{enumerate}
\end{enumerate}

\begin{enumerate}
\def\labelenumi{\alph{enumi})}
\setcounter{enumi}{1}
\item
  Let \(\mathbf{F}\) be a vector space, and let \(\mathbf{G}_1, \mathbf{G}_2\) be two vector subspaces of \(\mathbf{F}^*\) such that \(\mathbf{F}\) is in separating duality with \(\mathbf{G}_1\) and with \(\mathbf{G}_2\) . Show that, if \(\mathbf{F}\) is relatively bounded for \(\sigma(\mathbf{F}, \mathbf{G}_1)\) and \(\sigma(\mathbf{F}, \mathbf{G}_2)\) , it is so also for \(\sigma(\mathbf{F}, \mathbf{G}_1 + \mathbf{G}_2)\) .
\item
  Suppose that F has a countable basis. Show that, for every vector subspace G of \(\mathbf{F}^*\) which is in separating duality with F and has a countable basis, F is relatively bounded for \(\sigma(\mathrm{F},\mathrm{G})\) (by induction define two bases \((a_n)\) , \((b_n)\) of F and G respectively, such that \(\langle a_m, b_n \rangle = \delta_{mn}\) ).
\end{enumerate}

\begin{enumerate}
\def\labelenumi{\arabic{enumi})}
\setcounter{enumi}{2}
\tightlist
\item
  Let F be a vector space. Show that, for the topology \(\sigma(F, F^{*})\) , every bounded subset of F is finite dimensional. Deduce that, if F is infinite dimensional, there exist, in the completion \(\hat{F}\) of F (for \(\sigma(F, F^{*})\) ), compact subsets which are not contained in the closure of any bounded subset of F (cf.~II, p.~52, prop. 10).
\end{enumerate}

¶ 4) Let F and G be two vector spaces in separating duality, G (resp. F) being identified with the dual of F (resp. G) when the latter is assigned the topology σ(F, G) (resp. σ(G, F)). Let S (resp. T) be a covering of F (resp. G) consisting of convex, balanced and bounded subsets for σ(F, G) (resp. σ(G, F)). Show that the following propositions are equivalent :

\(\alpha)\) Every set \(\mathbf{M} \in \mathfrak{S}\) is precompact for the \(\mathfrak{T}\) -topology.

β) Every set \(\mathbf{N} \in \mathfrak{T}\) is precompact for the \(\mathfrak{S}\) -topology.

\(\gamma)\) On every set \(\mathbf{M} \in \mathfrak{S}\) , the topology induced by the \(\mathfrak{T}\) -topology is identical with the topology induced by \(\sigma(\mathrm{F}, \mathrm{G})\) .

\(\delta)\) On every set \(\mathbf{N} \in \mathfrak{T}\) , the topology induced by the \(\mathfrak{S}\) -topology is identical with the topology induced by \(\sigma(\mathbf{G}, \mathbf{F})\) .

(Use prop. 5 of III, p.~17 to show that \(\alpha\) ) implies \(\delta\) ) and exerc. 1 of II, p.~74 to show that \(\delta\) ) implies \(\beta\) .)

\begin{enumerate}
\def\labelenumi{\arabic{enumi})}
\setcounter{enumi}{4}
\item
  Let E and F be two locally convex Hausdorff spaces, S a family of subsets of E. In order that the S-topology on the space \(\mathcal{L}(\mathrm{E};\mathrm{F})\) be compatible with the vector space structure, it is necessary (and sufficient, cf.~III, p.~13, corollary) that every set of S is bounded in E.
\item
  \begin{enumerate}
  \def\labelenumii{\alph{enumii})}
  \tightlist
  \item
    Let E be a locally convex Hausdorff space. For every ultrafilter U on E, let \(U'\) be the filter for which the convex envelopes of sets of U form a base. Show that, in order that a point of E be a limit of U for the weakened topology, it is necessary and sufficient that it is a limit point of \(U'\) for the initial topology (use exerc. 11 of II, p.~84).
  \end{enumerate}
\end{enumerate}

\begin{enumerate}
\def\labelenumi{\alph{enumi})}
\setcounter{enumi}{1}
\tightlist
\item
  Let A be a convex subset of a vector space E and let \(T_{1}, T_{2}\) be two locally convex Hausdorff topologies on E, and \(T_{1}^{\prime}, T_{2}^{\prime}\) the corresponding weakened topologies. Show that if the topology induced on A by \(T_{1}\) is finer than the topology induced by \(T_{2}\) , then the topology induced on A by \(T_{1}^{\prime}\) is finer than the topology induced on A by \(T_{2}^{\prime}\) .
\end{enumerate}

¶ 7) Let \(\mathbf{F}\) be the direct sum space \(\mathbf{R}^{(\mathbf{N})}\) , \(\mathbf{G}\) the space \(\ell^1 (\mathbf{N})\) (I, p.~4, Example); \(\mathbf{F}\) and \(\mathbf{G}\) are put in duality by the bilinear form

\[
(x, y) \mapsto \sum_ {n} \xi_ {n} \eta_ {n}
\]

for \(x = (\xi_n)\in \mathbf{F}\) and \(y = (\eta_{n})\in \mathbf{G}\) .

\begin{enumerate}
\def\labelenumi{\alph{enumi})}
\tightlist
\item
  Show that, in F, every set K which is convex and compact for \(\sigma(F, G)\) is finite dimensional. (Assume the contrary; let \((n_k)\) be a strictly increasing sequence of integers \(>0\) , and \((a_k)\) a sequence of points of K such that the components of \(a_k\) for indices \(>n_k\) are zero, but that of index \(n_k\) is \(\neq 0\) . Show that there exists a sequence \((t_k)\) of numbers \(>0\) such that \(\sum_{k} t_k < +\infty\) and that in the Banach space \(\mathcal{B}(N)\) of all bounded sequences of real numbers,
\end{enumerate}

the point \(\sum_{k} t_k a_k\) has non-zero components for indices \(n_i\) , for all \(i\) ; deduce that the sequence

of partial sums \(s_{p} = \sum_{k=1}^{p} t_k a_k\) cannot have a limit point in F for \(\sigma(F, G)\) .

\begin{enumerate}
\def\labelenumi{\alph{enumi})}
\setcounter{enumi}{1}
\item
  Show that in F there exist compact sets for \(\sigma(\mathrm{F},\mathrm{G})\) which are infinite dimensional (observe that G is the dual of F for the topology induced by the topology of the normed space \(\mathcal{B}(\mathbf{N})\) ). Show that there exist precompact sets in F for \(\sigma(\mathrm{F},\mathrm{G})\) which are not relatively compact.
\item
  Deduce from a) and b) that \(\tau(\mathbf{G},\mathbf{F})=\sigma(\mathbf{G},\mathbf{F})\) , but that \(\tau(\mathbf{G},\mathbf{F})\) is distinct from the topology of uniform convergence on subsets of F which are compact for \(\sigma(\mathbf{F},\mathbf{G})\) .
\end{enumerate}

\begin{enumerate}
\def\labelenumi{\arabic{enumi})}
\setcounter{enumi}{7}
\tightlist
\item
  Let E be an infinite dimensional locally convex metrizable space, and \(E'\) its dual. In order that the topology \(\tau(\mathrm{E}, \mathrm{E}')\) be identical with the weakened topology \(\sigma(\mathrm{E}, \mathrm{E}')\) it is necessary and sufficient that E is isomorphic to an everywhere dense subspace of the product space \(R^{N}\) (resp. \(C^{N}\) ). (Observe that in \(E'\) , for the bornology consisting of equicontinuous sets, there exists a countable base (III, p.~1) consisting of finite dimensional sets; conclude that the vector space \(E'\) has a countable basis).
\end{enumerate}

Give an example of a locally convex Hausdorff space \(\mathbf{E}\) for which the initial topology, and the topologies \(\sigma(\mathrm{E},\mathrm{E}')\) and \(\tau(\mathrm{E},\mathrm{E}')\) are all distinct.

\begin{enumerate}
\def\labelenumi{\arabic{enumi})}
\setcounter{enumi}{8}
\tightlist
\item
  \begin{enumerate}
  \def\labelenumii{\alph{enumii})}
  \tightlist
  \item
    Let E be a locally convex, Hausdorff, bornological and quasi-complete space. In order that the topologies \(\tau(\mathrm{E}^{\prime},\mathrm{E})\) and \(\sigma(\mathrm{E}^{\prime},\mathrm{E})\) on the dual \(E^{\prime}\) of E be identical, it is necessary and sufficient that the topology of E is the finest locally convex topology (show that every bounded subset of E is finite dimensional).
  \end{enumerate}
\end{enumerate}

\begin{enumerate}
\def\labelenumi{\alph{enumi})}
\setcounter{enumi}{1}
\tightlist
\item
  Let E be a locally convex Hausdorff space, \(E'\) its dual. In order that the strong topology \(\beta(E', E)\) on \(E'\) be identical with the weak topology \(\sigma(E', E)\) , it is necessary and sufficient that the topology of the bornological space associated with E (III, p.~40, exerc. 1) is the finest locally convex topology on E.
\end{enumerate}

\begin{enumerate}
\def\labelenumi{\arabic{enumi})}
\setcounter{enumi}{9}
\tightlist
\item
  Let \(\mathbf{E}\) be a locally convex Hausdorff space, \(\mathbf{E}'\) its dual. Show that the following propositions are equivalent :
\end{enumerate}

α) E is barrelled;

β) every weakly bounded subset of E' is equicontinuous;

\(\gamma)\) every weakly bounded subset of \(\mathbf{E}'\) is relatively weakly compact, and the topology of \(\mathbf{E}\) is \(\tau (\mathbf{E},\mathbf{E}')\) .

\begin{enumerate}
\def\labelenumi{\arabic{enumi})}
\setcounter{enumi}{10}
\item
  Let E be a locally convex Hausdorff space, \(E'\) its dual, S a covering of E consisting of bounded subsets; we assign the S-topology to \(E'\) . Show that, for the bilinear form \((x, x') \mapsto \langle x, x' \rangle\) to be continuous on \(E \times E'\) , it is necessary and sufficient that the topology of E can be defined by a single norm, and that the S-topology is the strong topology on \(E'\) (cf.~III, p.~37, exerc. 2).
\item
  Let \(E\) be a locally convex Hausdorff space, \(E'\) its dual.
\end{enumerate}

\begin{enumerate}
\def\labelenumi{\alph{enumi})}
\item
  In order that there exists a weakly bounded and absorbent set in \(E'\) , it is necessary and sufficient that the topology of E is coarser than a normed space topology (cf.~IV, p.~47, exerc. 2).
\item
  In order that there exists an equicontinuous and weakly total set in \(E'\) , it is necessary and sufficient that the topology of E is finer than a normed space topology.
\item
  In order that there exists an equicontinuous absorbent set in \(E'\) , it is necessary and sufficient that the topology of E can be defined by a single norm.
\end{enumerate}

\begin{enumerate}
\def\labelenumi{\arabic{enumi})}
\setcounter{enumi}{12}
\tightlist
\item
  Let \(\mathbf{F}\) and \(\mathbf{G}\) be two vector spaces over \(\mathbf{R}\) in separating duality.
\end{enumerate}

\begin{enumerate}
\def\labelenumi{\alph{enumi})}
\item
  Let A be a convex set in F, not containing the origin, and compact for the topology \(\sigma(\mathbf{F}, \mathbf{G})\) ; let C be the convex cone with vertex 0 generated by A. Show that the polar cone \(C^{\circ}\) in G has an interior point for the topology \(\tau(\mathbf{G}, \mathbf{F})\) :
\item
  Conversely, let C be a proper convex cone with vertex 0, closed for \(\sigma(F, G)\) and having an interior point for the topology \(\tau(F, G)\) . Show that in G there exists a hyperplane H, closed for \(\sigma(G, F)\) , not containing the origin, such that \(H \cap C^{\circ}\) is compact for \(\sigma(G, F)\) and \((H \cap C^{\circ}) \cup \{0\}\) generates \(C^{\circ}\) .
\end{enumerate}

\begin{enumerate}
\def\labelenumi{\arabic{enumi})}
\setcounter{enumi}{13}
\tightlist
\item
  Let E be a locally convex Hausdorff and quasi-complete space, \(E'\) its dual. Show that on \(E_{1}'\) the topology of compact convergence is the finest of the topologies compatible with the duality between E and \(E'\) and which induces the same topology as \(\sigma(E', E)\) on every equicontinuous subset of \(E'\) (cf.~IV, p.~48, exerc. 4).
\end{enumerate}

¶ 15) a) Let E be a locally convex Hausdorff space, A a convex and balanced set in E, and u a linear form on E. Show that if \(\mathrm{A} \cap u^{-1}(0)\) is closed with respect to A, then the restriction of u to A is continuous. (If not, show that 0 will be in the closure of the intersection of A and a hyperplane \(u^{-1}(\alpha)\) with \(\alpha \neq 0\) ; deduce that if \(b \in A\) is such that \(u(b) = -\alpha\) , the point \(\frac{1}{2}b\) will be in the closure of \(\mathrm{A} \cap u^{-1}(0)\) .)

\begin{enumerate}
\def\labelenumi{\alph{enumi})}
\setcounter{enumi}{1}
\item
  Let E be a real locally convex Hausdorff and complete space, \(E'\) its dual. Show that if a hyperplane \(H'\) of \(E'\) is such that its intersection with every equicontinuous and weakly closed subset \(M' \subset E'\) is weakly closed, then \(H'\) is weakly closed (use a) and III, p.~21, cor. 1).
\item
  Let E be a locally convex Hausdorff space, \(E'\) its dual and C a convex, balanced and closed set in E. Let u be a linear form on E; show that if the restriction of u to C is continuous for the initial topology, it is also continuous for \(\sigma(\mathrm{E}, \mathrm{E}')\) (use a)). Show by an example that the restriction of u to the vector subspace M generated by C is not necessarily continuous (take \(E = R^{(N)}\) with the norm \(\|x\| = \sup_{n} |\xi_{n}|\) , and for C take a suitable convex set generating E).
\end{enumerate}

\begin{enumerate}
\def\labelenumi{\arabic{enumi})}
\setcounter{enumi}{15}
\tightlist
\item
  Let F and G be two vector spaces in separating duality, the set of all linear forms \(x'\) on F which are bounded on every subset of F bounded for \(\sigma(F, G)\) is called the enclosure of G in the algebraic dual \(F^{*}\) of F; this is a vector subspace \(\tilde{G}\) of \(F^{*}\) , which is the dual of F when F is assigned the topology of the bornological space associated (III, p.~40, exerc. 1) with one of the topologies compatible with the duality between F and G. Then G is said to be enclosed in \(F^{*}\) if \(\tilde{G} = G\) .
\end{enumerate}

\begin{enumerate}
\def\labelenumi{\alph{enumi})}
\tightlist
\item
  Let M be a closed vector subspace of F for the topology \(\sigma(\mathrm{F},\mathrm{G})\) . Show that if G is enclosed in \(F^{*}\) , and if F/M is assigned the topology \(\sigma(\mathrm{F}/\mathrm{M},\mathrm{M}^{\circ})\) , then its dual is enclosed in \((\mathrm{F}/\mathrm{M})^{*}\) .
\item
  Let E be a Hausdorff locally convex space, \(E'\) its dual. For E to be bornological, it is necessary and sufficient that its topology is identical with \(\tau(\mathrm{E},\mathrm{E}')\) and that \(E'\) is enclosed in \(E^{*}\) .
\item
  Let \((\mathrm{E}_{i})_{i\in\mathbb{I}}\) be a family of Hausdorff locally convex spaces, and F the direct sum space of the \(E_{i}\) , endowed with the topology defined in I, p.~24, exerc. 14. Show that the dual of F is canonically isomorphic to the subspace of the product \(\prod_{i\in I}E_{i}'\) of the duals of the \(E_{i}\) , consisting
\end{enumerate}

of all families \((x_{i}^{\prime})\) such that \(x_{i}^{\prime}=0\) except for a countable number of indices. (Let \(V_{i}\) be an arbitrary neighbourhood of 0 in \(E_{i}\) . Show that if \(x_{i}^{\prime}\neq0\) for an uncountable set of indices, then there exists a number \(\alpha>0\) and an uncountable set H⊂I such that there exists \(x_{i}\in V_{i}\) for which \(\langle x_{i},x_{i}^{\prime}\rangle\geqslant\alpha\) for all \(i\in H\) ; conclude that \((x_{i}^{\prime})\) cannot belong to F'.)

\begin{enumerate}
\def\labelenumi{\alph{enumi})}
\setcounter{enumi}{3}
\tightlist
\item
  Show that if I is uncountable then \(\mathbf{F}'\) is not enclosed in \(\mathbf{F}^*\) ; if we take \(\mathrm{E}_i = \mathbf{R}\) for all \(i \in \mathrm{I}\) , then \(\mathbf{F}'\) endowed with the strong topology, is not complete, and there exist strongly bounded subsets in \(\mathbf{F}'\) which are not weakly relatively compact.
\end{enumerate}

\begin{enumerate}
\def\labelenumi{\arabic{enumi})}
\setcounter{enumi}{16}
\item
  Let E be a Hausdorff locally convex space, \(E'\) its dual and, in \(E'\) , let \(B_{1}\) be the family of convex equicontinuous sets, \(B_{2}\) the family of convex, relatively weakly compact subsets, \(B_{3}\) the family of convex strongly bounded subsets, and \(B_{4}\) the family of convex weakly bounded subsets. Then \(B_{1} \subset B_{2} \subset B_{3} \subset B_{4}\) . Give an example of a space E for which the four families of sets are distinct (take for E a product of three spaces for which we have, respectively \(B_{1} \neq B_{2}\) (cf.~IV, p.~49, exerc. 8), \(B_{2} \neq B_{3}\) (exerc. 16, d)), \(B_{3} \neq B_{4}\) (III, p.~23, Remark 2).
\item
  Let E, F be two vector spaces and \(E^{*}\) , \(F^{*}\) their respective algebraic duals. Show that if u is a linear mapping from \(F^{*}\) into E, which is continuous for the topologies \(\sigma(F^{*}, F)\) and \(\sigma(E, E^{*})\) , then \(u(F^{*})\) is finite dimensional. (Use prop. 2 of IV, p.~27 to show that u is a strict morphism, and deduce that \(u(F^{*})\) is a subspace of E of minimal type (II, p.~85, exerc. 13); conclude by considering the bounded sets of this subspace.)
\item
  Let E and F be two Hausdorff locally convex spaces, \(E'\) and \(F'\) their duals. For every subset H of the space \(\mathcal{L}(\mathrm{E};\mathrm{F})\) of continuous linear mappings from E into F, let \({}^{t}H\) denote the set of the transposes of the mapping \(u\in H\) . For every subset M (resp. \(N'\) ) of E (resp. \(F'\) ), let \(\mathrm{H}(\mathrm{M})\) (resp. \({}^{t}\mathrm{H}(\mathrm{N}')\) ) denote the union of the sets \(u(\mathrm{M})\) (resp. \({}^{t}u(\mathrm{N}')\) ) as u ranges over H. a) For H to be equicontinuous, it is necessary and sufficient that, for every equicontinuous subset \(N'\) of \(F'\) , \({}^{t}\mathrm{H}(\mathrm{N}')\) is an equicontinuous subset of \(E'\) .
\end{enumerate}

\begin{enumerate}
\def\labelenumi{\alph{enumi})}
\setcounter{enumi}{1}
\item
  Let \(\mathfrak{S}\) be a set of bounded subsets of E. Show that H is bounded in \(\mathcal{L}(\mathrm{E};\mathrm{F})\) for the \(\mathfrak{S}\) -topology, if and only if for every \(y' \in \mathrm{F}'\) , \({}^{t}\mathrm{H}(y')\) is bounded in \(\mathrm{E}'\) for the \(\mathfrak{S}\) -topology.
\item
  Let S be a set of bounded subsets of E, J a set of bounded subsets of F forming an adapted bornology of F (III, p.~3, def. 4). Suppose E' is assigned the S-topology and F' the J-topology. Then \({}^{t}\mathrm{H}\) is equicontinuous if and only if for every set B ∈ S, H(B) belongs to J. In particular, \({}^{t}\mathrm{H}\) is equicontinuous for the strong topologies on F' and E' if and only if H is bounded in L(E; F) for the topology of bounded convergence.
\end{enumerate}

\(d\) ) Deduce from \(b\) ) and \(c\) ) that, if \({}^t\mathrm{H}\) is bounded for the topology of simple convergence in \(\mathcal{L}(\mathrm{F}';\mathrm{E}')\) when \(\mathrm{F}'\) and \(\mathrm{E}'\) are assigned the strong topologies, then \({}^t\mathrm{H}\) is equicontinuous for these topologies.

\begin{enumerate}
\def\labelenumi{\alph{enumi})}
\setcounter{enumi}{4}
\tightlist
\item
  Show that if E is barrelled, the following properties are equivalent :
\end{enumerate}

α) H is simply bounded in \(\mathcal{L}(E; F)\);

β) H is equicontinuous;

\(\gamma)\) \({}^{t}\mathrm{H}\) is simply bounded in \(\mathcal{L}(\mathrm{F}';\mathrm{E}')\) when \(\mathrm{E}'\) is assigned the weak topology \(\sigma (\mathrm{E}',\mathrm{E})\)

\(\delta)\) \({}^{t}\mathrm{H}\) is equicontinuous when \(\mathrm{E}'\) and \(\mathrm{F}'\) are assigned the strong topologies.

\begin{enumerate}
\def\labelenumi{\alph{enumi})}
\setcounter{enumi}{5}
\tightlist
\item
  Show that if E is infrabarrelled (III, p.~44, exerc. 7), then the properties \(\beta\) and \(\delta\) of \(e\) ) are equivalent, and are also equivalent to the following :
\end{enumerate}

ε) H is bounded in \(\mathcal{L}(\mathrm{E};\mathrm{F})\) for the topology of bounded convergence;

\(\phi)\) \({}^{t}\mathrm{H}\) is simply bounded in \(\mathcal{L}(\mathrm{F}';\mathrm{E}')\) when \(\mathrm{E}'\) is assigned the strong topology.

\begin{enumerate}
\def\labelenumi{\alph{enumi})}
\setcounter{enumi}{6}
\tightlist
\item
  Show that if E is quasi-complete, the properties \(\alpha\) ), \(\gamma\) ) and \(\delta\) ) of e) are equivalent.
\end{enumerate}

\begin{enumerate}
\def\labelenumi{\arabic{enumi})}
\setcounter{enumi}{19}
\tightlist
\item
  Let \(\mathbf{E}\) be a Hausdorff locally convex space, and \(\mathbf{E}'\) its dual.
\end{enumerate}

\begin{enumerate}
\def\labelenumi{\alph{enumi})}
\item
  Let M be a closed vector subspace of E. The topology \(\tau(\mathbf{M}, \mathbf{E}'/\mathbf{M}^{\circ})\) is identical with the topology induced on M by \(\tau(\mathbf{E}, \mathbf{E}')\) if and only if every convex balanced set in \(E'/M^{\circ}\) which is compact for the weak topology \(\sigma(\mathrm{E}'/\mathrm{M}^{\circ}, \mathrm{M})\) , is the canonical image of a convex balanced, compact (for \(\sigma(\mathrm{E}', \mathrm{E})\) ) subset of \(E'\) (cf.~V, p.~73, exerc. 15).
\item
  Let N be a dense vector subspace of E. If the topology induced on N by that of E is identical with \(\tau(\mathrm{N}, \mathrm{E}')\) , show that the topology of E is identical with \(\tau(\mathrm{E}, \mathrm{E}')\) .
\end{enumerate}

\begin{enumerate}
\def\labelenumi{\arabic{enumi})}
\setcounter{enumi}{20}
\tightlist
\item
  \begin{enumerate}
  \def\labelenumii{\alph{enumii})}
  \tightlist
  \item
    Let E be a vector space, \(E^{*}\) its algebraic dual. Show that the topology \(\tau(\mathrm{E},\mathrm{E}^{*})\) is the finest locally convex topology and that the topology \(\tau(\mathrm{E}^{*},\mathrm{E})\) is identical with \(\sigma(\mathrm{E}^{*},\mathrm{E})\) . b) Let \(E^{**}\) be the algebraic dual of \(E^{*}\) . Show that if E is infinite dimensional, then E is dense in \(E^{**}\) for all the topologies compatible with the duality between \(E^{**}\) and \(E^{*}\) , but that the topology induced on E by \(\tau(\mathrm{E}^{**},\mathrm{E}^{*})\) is distinct from \(\tau(\mathrm{E},\mathrm{E}^{*})\) .
  \end{enumerate}
\item
  Let E be a Hausdorff locally convex space, \(E'\) its dual and M a closed subspace of E. a) Show that if the closed convex envelope of a compact set in E is compact, then the topology of compact convergence on \(E'/M^{\circ}\) (identified with the dual of M) is the quotient topology of the topology of compact convergence on \(E'\) by \(M^{\circ}\) .
\end{enumerate}

\begin{enumerate}
\def\labelenumi{\alph{enumi})}
\setcounter{enumi}{1}
\tightlist
\item
  Show that if \(\mathbf{M}\) is infrabarrelled and if \(\mathrm{E}^{\prime} / \mathrm{M}^{\circ}\) , endowed with the topology \(\beta (\mathrm{E}^{\prime} / \mathrm{M}^{\circ},\mathbf{M})\) is bornological, then the topology \(\beta (\mathrm{E}^{\prime} / \mathrm{M}^{\circ},\mathbf{M})\) is the quotient of \(\beta (\mathrm{E}^{\prime},\mathrm{E})\) by \(\mathbf{M}^{\circ}\) .
\end{enumerate}

\begin{enumerate}
\def\labelenumi{\arabic{enumi})}
\setcounter{enumi}{22}
\item
  Give an example of a family \((\mathrm{E}_{i})_{i\in\mathrm{I}}\) of Hausdorff locally convex spaces such that the canonical mapping from \(\bigoplus_{i\in I} E_{i}^{\prime}\) onto the dual of \(P = \prod_{i\in I} E_{i}\) is not an isomorphism from the topological direct sum of the \(E_{i}^{\prime}\) endowed with the weak topologies \(\sigma(\mathrm{E}_{i}^{\prime}, \mathrm{E}_{i})\) onto the dual \(P^{\prime}\) endowed with \(\sigma(\mathrm{P}^{\prime}, \mathrm{P})\) .
\item
  Let \((\mathrm{E}_{\alpha})_{\alpha\in\mathrm{A}}\) be a family of locally convex Hausdorff spaces, E a vector space and for every \(\alpha\in A\) , let \(f_{\alpha}\) be a linear mapping from \(E_{\alpha}\) into E. On E consider the finest locally convex topology T for which the functions \(f_{\alpha}\) are continuous (II, p.~27); suppose T is Hausdorff and let \(E_{\alpha}^{\prime}\) denote the dual of \(E_{\alpha}\) , and \(E^{\prime}\) that of E endowed with T. Show that if, for every \(\alpha\in A\) , the topology of \(E_{\alpha}\) is identical with \(\tau(\mathrm{E}_{\alpha},\mathrm{E}_{\alpha}^{\prime})\) , then the topology T is identical with \(\tau(\mathrm{E},\mathrm{E}^{\prime})\) .
\item
  \begin{enumerate}
  \def\labelenumii{\alph{enumii})}
  \tightlist
  \item
    Show that every real (resp. complex) Banach space is isometric to a closed vector subspace of a Banach space of the form \(\mathcal{C}(\mathbf{S};\mathbf{R})\) (resp. \(\mathcal{C}(\mathbf{S};\mathbf{C})\) ) consisting of all real (resp. complex) continuous functions defined on a compact space S (GT, X, § 4, No.~1) (use formula (3) of IV, p.~7).
  \end{enumerate}
\end{enumerate}

\begin{enumerate}
\def\labelenumi{\alph{enumi})}
\setcounter{enumi}{1}
\tightlist
\item
  Deduce from a) that every Hausdorff locally convex space E is isomorphic to a subspace of a locally convex space of the form \(\mathcal{C}_{c}(\mathrm{L};\mathbf{R})\) (resp. \(\mathcal{C}_{c}(\mathrm{L};\mathbf{C})\) ) (GT, X, § 1, No.~6). In particular, every Fréchet space is isomorphic to a closed subspace of a space \(\mathcal{C}_{c}(\mathrm{L};\mathbf{R})\) (resp. \(\mathcal{C}_{c}(\mathrm{L};\mathbf{C})\) ), where L is locally compact and separable.
\end{enumerate}
% semantic-fragments: legacy-input-seam
\relax{}
